\documentclass[11pt]{article}
\usepackage[margin=25mm]{geometry}
\usepackage{fontspec}
\setmainfont{Latin Modern Roman}
\usepackage{amsmath,amssymb,mathtools}
\usepackage{unicode-math}
\setmathfont{Latin Modern Math}
\usepackage{xurl,hyperref}
\hypersetup{colorlinks=true,urlcolor=blue}
\setcounter{secnumdepth}{0}
\setlength{\parindent}{0pt}
\setlength{\parskip}{5pt}
\setlength{\emergencystretch}{3em}
\newcommand{\sref}[2]{\hyperlink{source:#1:#2}{[S#2]}}
\newcommand{\cataloguescope}[1]{\paragraph{Recorded target.}#1}
\newcommand{\R}{\mathbb R}
\newcommand{\C}{\mathbb C}
\newcommand{\N}{\mathbb N}
\newcommand{\Z}{\mathbb Z}
\newcommand{\Q}{\mathbb Q}
\newcommand{\F}{\mathbb F}
\newcommand{\eps}{\varepsilon}
\begin{document}
\section{Formal, Jacobi, and wave-function transformations of the string coupling}\label{statement}

\subsection{Definitions and mathematical statement}

\paragraph{Shared notation.}$\mathbb N=\{1,2,\ldots\}$; the standard number fields have their usual meanings. All additional conventions are stated locally.


\paragraph{Definitions and notation.}The source’s $F(\tau,\lambda)=\sum_g f_g(\tau)\lambda^{2g-2}$ is a genus expansion in a string-coupling parameter $\lambda$; $\tau$ is a modular parameter. A formal expansion need not define a convergent function. A transformation law must specify the operation on $\lambda$, its simultaneous action on $\tau$ and the multiplier or weight.

\paragraph{Catalogue wording retained for traceability.}
Does $F= \sum f_g(\tau)\lambda^{2g-2}$ have transformation properties with respect to $\lambda$?
\paragraph{Scope and source qualification.}The source does not identify the proposed transformations or assert convergence of the genus series. Formal identities, Jacobi transformation laws and wave-function interpretations must be kept distinct. The catalogue title/status was paired with a machine-generated research attempt; that attempt is not peer-reviewed evidence of settlement. The statement covers the original source questions rather than only the title’s subcase.
\subsection{Short English statement}
Determine whether the full genus expansion has meaningful transformation laws involving the string coupling.
\subsection{Research attempt}
\paragraph{Model, date and investigated scope.}
OpenAI GPT-6 Astra (configured model \texttt{gpt-6-astra}), ultra reasoning; generated 11 October 2026. Proved an exact equivalence between coefficient modular weights and a formal lambda-scaling law, and identified the obstruction to formal elliptic translations. No target coefficient family or analytic transformation is established.

\paragraph{Sources and literature check.}
Read the supplied complete definition, catalogue wording, scope qualifications and source list for this canonical ID. The original-source citations are inherited as supplied references; no fresh primary-source or novelty search was performed for this attempt. No current open-status claim is inferred from this calculation.

\paragraph{A formal transformation criterion.}Treat \(F(\tau,\lambda)=\sum_{g\ge0}f_g(\tau)\lambda^{2g-2}\) as a formal Laurent series with only the possible order-two pole. Let \(\gamma\tau=(a\tau+b)/(c\tau+d)\), and suppose, as an explicit hypothesis, that each coefficient transforms by
\[f_g(\gamma\tau)=(c\tau+d)^{k+2g-2}f_g(\tau).
\]
Substituting coefficient by coefficient proves the formal identity
\[F(\gamma\tau,\lambda/(c\tau+d))=(c\tau+d)^kF(\tau,\lambda).
\]
Conversely, equality of these formal Laurent series forces exactly the individual coefficient laws. Thus this scaling transformation neither adds nor removes modular information: it packages the specified weights consistently. Also \(F(\tau,-\lambda)=F(\tau,\lambda)\) is automatic from the even powers, with no modular hypothesis.
\paragraph{Why this is not a Jacobi law.}A genuine Jacobi transformation can contain an exponential quadratic factor in \(\lambda\) and elliptic translations \(\lambda\mapsto\lambda+m\tau+n\). Such translations are not automatically defined on an arbitrary infinite formal series about zero: a fixed low-order coefficient after a nonzero shift may receive infinitely many contributions. Nor does the weight hypothesis allow the quadratic multiplier unless corresponding mixing among the \(f_g\) is present. Convergence or a specified formal completion is needed.
\paragraph{Gap.}The source specifies neither the actual coefficient family nor a transformation group. Proving their coefficient transformation laws, or deriving the anomaly-driven mixing and analytic meaning required for a wave-function/Jacobi interpretation, is the first missing step. The algebra above is a concrete conditional test, not evidence that an unspecified genus expansion satisfies such a law.

\paragraph{Verification and adversarial audit.}
Checked genus-zero Laurent pole, weight powers, converse coefficient comparison, parity and infinite-contribution obstruction under translation. This was an LLM adversarial audit, not independent human verification; no theorem prover was used. Prompts A then B applied to this ID: Checked genus-zero Laurent pole, weight powers, converse coefficient comparison, parity and infinite-contribution obstruction under translation. Scope remains limited: Determine and prove the actual coefficient laws or anomaly mixing and specify a convergent or completed setting for nonzero lambda translations.

Coordinator checked the formal coefficient transformation equivalence and the obstruction to unqualified nonzero elliptic translations of a formal expansion. This additional check was model-assisted, not independent human verification.

\paragraph{Tools, process and contributions.}
Python 3 and local shell used to read supplied files and serialize this manuscript; mathematics derived symbolically in the text, with no numerical experiment or formal proof assistant. Contribution prompt A was followed by prompt B. GPT-6 Astra with ultra reasoning drafted the research and performed its stated self-audit; any additional coordinator contribution is identified in the verification and audit. A GPT-6-based Codex coordinator assembled the manuscript and checked repository integration; its exact variant and reasoning configuration were not exposed. No human mathematical contribution, independent human proof check or formal proof certification is supplied. Complete derivations and the consulted source references are included here; private reasoning transcripts are not published.

\paragraph{Final status and first remaining gap.}
\textbf{UNRESOLVED}. Determine and prove the actual coefficient laws or anomaly mixing and specify a convergent or completed setting for nonzero lambda translations.

\section{Completion Estimate}
\noindent\textbf{COMPLETION: 0\%}
The calculation is an elementary benchmark or conditional reduction; it establishes no claimed new portion of the complete source target.

\subsection{Sources}
\begin{itemize}
\item[\hypertarget{source:statement:1}{S1}] ULAM.AI and the UnsolvedMath Contributors, UnsolvedMath, catalogue entry. Catalogue curation under CC BY 4.0; definitions and English exposition added locally.
\url{https://huggingface.co/datasets/ulamai/UnsolvedMath}.
\item[\hypertarget{source:statement:2}{S2}] American Institute of Mathematics, Gromov-Witten invariants and number theory, source item 4.3.
\url{http://aimpl.org/gromwitnumthry/4/}.
\end{itemize}
\label{end:statement}
\end{document}
